\documentclass[12pt,oneside]{article}
\usepackage[utf-8]{inputenc}
\usepackage[spanish]{babel}
\usepackage[margin=1.5cm]{geometry}
\usepackage{hyperref}
\usepackage{fancyhdr}
\usepackage{xcolor}
\usepackage{booktabs}
\usepackage{tcolorbox}
\usepackage{array}
\usepackage{makecell}
\usepackage{listings}
\usepackage{enumitem}

% Colores
\definecolor{darkblue}{RGB}{0,51,102}
\definecolor{lightblue}{RGB}{230,240,250}
\definecolor{accent}{RGB}{255,102,0}
\definecolor{darkgreen}{RGB}{0,100,0}
\definecolor{lightgreen}{RGB}{200,255,200}

% Encabezado y pie
\pagestyle{fancy}
\fancyhf{}
\fancyhead[L]{\small TMP AI Assistant v3.0}
\fancyhead[R]{\small\textcolor{darkblue}{Manual Usuario Final}}
\fancyfoot[C]{\small\thepage}
\renewcommand{\headrulewidth}{0.5pt}
\renewcommand{\footrulewidth}{0.5pt}

% Títulos
\usepackage{titlesec}
\titleformat{\section}{\LARGE\bfseries\color{darkblue}}{}{0em}{\colorbox{lightblue}{\parbox{\dimexpr\textwidth-2\fboxsep}{\ \thesection.\ #1\ }}}
\titleformat{\subsection}{\Large\bfseries\color{darkblue}}{\thesubsection.\ }{0em}{}
\titleformat{\subsubsection}{\large\bfseries\color{accent}}{\thesubsubsection.\ }{0em}{}

\usepackage{tocbibind}
\setcounter{tocdepth}{2}

\title{
    \Huge\bfseries\textcolor{darkblue}{TMP AI ASSISTANT}\\
    \Large\textcolor{accent}{Manual de Usuario Final Completo}\\
    \vspace{0.5cm}
    \large\textit{Versión 3.0 | Enero 2026 | EDICIÓN COMPLETA}
}
\author{\textcolor{darkblue}{Desarrollado por Alfredo}}
\date{\today}

\begin{document}

\maketitle
\thispagestyle{empty}

\begin{center}
    \vspace{1cm}
    \Large\textbf{\textcolor{darkblue}{Automatiza tu contabilidad en QuickBooks Online}}
    \vspace{0.5cm}
    
    \large\textcolor{darkblue}{Guía completa con flujos de trabajo, FAQ y troubleshooting avanzado}
    
    \vspace{2cm}
    \begin{tcolorbox}[colback=lightblue, colframe=darkblue, boxrule=2pt, width=0.95\textwidth]
        \textbf{\Large\textcolor{darkblue}{🚀 VERSIÓN COMPLETA}}
        
        Manual expandido con 15+ secciones, flujos de trabajo paso a paso, FAQ exhaustivo, y troubleshooting para 20+ errores comunes.
    \end{tcolorbox}
\end{center}

\newpage
\tableofcontents
\newpage

% ============================================
% SECCIÓN 1: OVERVIEW
% ============================================
\section{Overview del Proyecto}

\subsection{¿Qué es TMP AI Assistant?}

TMP AI Assistant es un asistente contable conversacional que automatiza operaciones en QuickBooks Online mediante lenguaje natural en español, integrando:

\begin{itemize}
    \item ✅ 31 funciones inteligentes (búsquedas, transacciones, reportes, OCR)
    \item ✅ Integración directa con QB API (OAuth 2.0)
    \item ✅ Procesamiento batch de CSV (100+ transacciones/segundo)
    \item ✅ OCR inteligente de facturas PDF
    \item ✅ Optimización de tokens (57\% menos costo)
    \item ✅ Reconciliación bancaria automática
    \item ✅ Reportes financieros en devengado y caja
\end{itemize}

\subsection{Ventajas competitivas principales}

\begin{table}[h]
\centering
\begin{tabular}{|p{3cm}|p{9cm}|}
\hline
\textbf{Característica} & \textbf{Ventaja} \\
\hline
Lenguaje Natural Español & Interfaz conversacional sin curva de aprendizaje \\
\hline
API Integration & Datos sincronizados en tiempo real (sin doble entrada) \\
\hline
Batch Processing & 500 transacciones en 30 segundos (vs manual: 8 horas) \\
\hline
OCR Inteligente & Escanea PDFs y genera transacciones automáticas \\
\hline
Costo Optimizado & \$10 dura 6.3 meses (vs 2.7 sin optimizar) \\
\hline
Reconciliación Auto & Valida bancos en minutos (vs 2 horas manual) \\
\hline
\end{tabular}
\caption{Ventajas vs métodos tradicionales}
\end{table}

\newpage

% ============================================
% SECCIÓN 2: INSTALACIÓN Y CONFIGURACIÓN
% ============================================
\section{Instalación y Configuración Completa}

\subsection{Requisitos previos}

\begin{tcolorbox}[colback=yellow!10, colframe=accent, boxrule=1.5pt]
    \begin{enumerate}[leftmargin=*]
        \item Python 3.9 o superior (\texttt{python3 --version})
        \item QuickBooks Online (no Desktop)
        \item OAuth 2.0 credentials de QB
        \item API key de OpenRouter
        \item Conexión a internet estable (>1 Mbps)
        \item ~100 MB de espacio en disco
    \end{enumerate}
\end{tcolorbox}

\subsection{Paso 1: Obtener credenciales de QuickBooks}

\begin{enumerate}
    \item Ve a \textbf{developer.intuit.com}
    \item Crea una aplicación OAuth 2.0
    \item Copia los valores:
    \begin{itemize}
        \item Client ID
        \item Client Secret
        \item Realm ID (tu empresa)
    \end{itemize}
    \item Autentica en QB Online para obtener access\_token y refresh\_token
\end{enumerate}

\subsection{Paso 2: Obtener API key de OpenRouter}

\begin{enumerate}
    \item Ve a \textbf{openrouter.ai}
    \item Crea cuenta (o inicia sesión)
    \item Genera API key en settings
    \item Copia el valor \texttt{sk-or-v1-xxxxx}
\end{enumerate}

\subsection{Paso 3: Crear archivo .env}

En la carpeta raíz del proyecto, crea archivo \textbf{.env}:

\begin{verbatim}
QB_ACCESS_TOKEN=eyJkb3JkYWduIjp...
QB_REFRESH_TOKEN=L0L507KR...
QB_CLIENT_ID=ABCDxxxxEFGH
QB_CLIENT_SECRET=xxxxyyyy1234
QB_REALM_ID=1234567890
OPENROUTER_API_KEY=sk-or-v1-xxxxx
QB_ENVIRONMENT=production
LOG_LEVEL=INFO
\end{verbatim}

\subsection{Paso 4: Instalar dependencias}

\begin{verbatim}
pip install -r requirements.txt
\end{verbatim}

Librerías necesarias:
\begin{itemize}
    \item requests (HTTP)
    \item python-dotenv (configuración)
    \item pandas (CSV processing)
    \item openpyxl (Excel export)
    \item PyPDF2 (OCR support)
\end{itemize}

\subsection{Paso 5: Crear estructura de carpetas}

\begin{verbatim}
proyecto/
├── main.py
├── ocr_bills.py
├── .env (NO compartir)
├── requirements.txt
├── CONTEXT.md
├── Pending bills/      (PDFs sin procesar)
├── Processed bills/    (PDFs ya procesados)
├── Backups/           (respaldos)
└── outputs/           (reportes generados)
\end{verbatim}

\subsection{Paso 6: Verificar instalación}

\begin{verbatim}
python3 main.py

Esperado:
✅ Conectado a QuickBooks Online
✅ Chart of Accounts cargado (N cuentas)
✅ Sesión iniciada
\end{verbatim}

\newpage

% ============================================
% SECCIÓN 3: OPERACIONES PRINCIPALES
% ============================================
\section{Operaciones Principales (Guía Práctica)}

\subsection{1. Búsquedas Inteligentes}

\subsubsection{Comando}
\begin{verbatim}
Busca [cliente|proveedor|cuenta|item] [nombre parcial]
\end{verbatim}

\subsubsection{Ejemplo avanzado}

\begin{verbatim}
Usuario: "Busca todos los clientes que empiecen con ACME"

Respuesta:
✅ Encontrados 5 resultados:
1. ACME Corp (ID: 123) - Cliente activo
2. ACME International (ID: 456) - Cliente
3. ACME Services (ID: 789) - Proveedor
4. Acme Mexico (ID: 101) - Suspendido
5. AcmePlus (ID: 102) - Cliente

Información: Usa fuzzy matching (búsqueda aproximada).
El sistema entiende:
- Variaciones de mayúscula/minúscula
- Errores de digitación (acme ≈ acmi)
- Nombres parciales (acm = acme, acme corp, etc.)
\end{verbatim}

\subsubsection{Casos de uso}

\begin{itemize}
    \item Verificar si cliente existe antes de crear factura
    \item Obtener ID exacto de proveedor
    \item Buscar cuentas por categoría (\textbf{búsca cuenta equity})
    \item Encontrar items por descripción
\end{itemize}

\subsection{2. Crear Facturas (Invoices)}

\subsubsection{Variantes de comando}

\begin{enumerate}
    \item \textbf{Monto fijo:} \texttt{Crea factura a [Cliente] por \$[Monto] por [Concepto]}
    \item \textbf{Por horas:} \texttt{Haz factura a [Cliente] por [Horas] hrs a \$[Precio/hr]}
    \item \textbf{Especificando fecha:} \texttt{Crea factura a [Cliente] por \$[Monto], fecha [DD/MM]}
    \item \textbf{Vencimiento personalizado:} \texttt{Factura a [Cliente] \$[Monto], vence en [Días] días}
\end{enumerate}

\subsubsection{Ejemplo completo: Factura de servicios}

\begin{verbatim}
Usuario: "Crea factura a ACME Corp por 3000, concepto 
desarrollo web, vence 45 días, para el 20 de enero"

Asistente:
🔍 Verificando datos...
✅ Cliente encontrado: ACME Corp (ID: 123)
✅ Cuenta de ingreso: Software Development Services
✅ Monto: $3,000.00

📋 Resumen:
   Cliente: ACME Corp
   Concepto: Desarrollo web
   Monto: $3,000.00
   Fecha: 20/01/2026
   Vencimiento: 04/03/2026 (45 días)
   Referencia: INV-1045

¿Confirmas? (sí/no)
\end{verbatim}

\subsubsection{Buenas prácticas}

\begin{itemize}
    \item ✅ Menciona concepto específico (no solo \textbf{servicio})
    \item ✅ Revisa monto antes de confirmar
    \item ✅ Usa fechas en formato DD/MM
    \item ❌ No crees facturas con impuesto (se agrega después en QB)
\end{itemize}

\subsection{3. Crear Bills (Cuentas por Pagar)}

\subsubsection{Ejemplo avanzado: Proveedor recurrente}

\begin{verbatim}
Usuario: "Registra bill a Telecom Venezuela por 1500 
por servicio telefónico enero, vencimiento 10 febrero"

Asistente:
✅ Bill creado:
   Proveedor: Telecom Venezuela
   Monto: $1,500.00
   Categoría de gasto: Utilities
   Vencimiento: 10/02/2026
   Referencia: BILL-2345
\end{verbatim}

\subsubsection{Casos de uso especiales}

\begin{enumerate}
    \item \textbf{Servicio recurrente:} \texttt{Bill a [Proveedor] \$[Monto] mensual}
    \item \textbf{Gastos múltiples líneas:} Procesar con CSV en su lugar
    \item \textbf{Pago anticipado:} Bill con vencimiento futuro
\end{enumerate}

\subsection{4. Depósitos: Guía completa}

\subsubsection{Tipos de depósito soportados}

\begin{table}[h]
\centering
\begin{tabular}{|l|p{7cm}|}
\hline
\textbf{Tipo} & \textbf{Cuándo usar} \\
\hline
Depósito simple & 1-2 clientes pagando \\
\hline
Batch CSV & 3+ clientes, múltiples depósitos \\
\hline
Bank Feed splits & Stripe/PayPal con múltiples ventas + fees \\
\hline
Transferencia interna & Mover dinero entre cuentas de la empresa \\
\hline
\end{tabular}
\caption{Tipos de depósitos}
\end{table}

\subsubsection{Flujo completo: Batch de 50 depósitos}

\begin{enumerate}
    \item Ejecutar: \texttt{template csv}
    \item Generar plantilla: \texttt{depositstemplate.csv}
    \item Llenar en Excel (50 filas de clientes)
    \item Validar:
    \begin{itemize}
        \item Nombres de clientes exactos
        \item Montos con 2 decimales
        \item Fechas en formato YYYY-MM-DD
        \item Cuentas origen/destino existen
    \end{itemize}
    \item Ejecutar: \texttt{Procesa deposits\_enero.csv}
    \item Confirmación (revisa totales)
    \item Sistema crea 50 depósitos en QB
    \item Archivo procesado se guarda con timestamp
\end{enumerate}

\subsubsection{Ejemplo CSV real}

\begin{verbatim}
customer_name,amount,from_account,to_account,date,memo
Cliente A,1500.50,Client Retainers,Checking,2026-01-22,Pago enero
Cliente B,2300.00,Prepaid Labour,Checking,2026-01-22,Anticipo
Freelancer X,500.00,1099 Retainers,Checking,2026-01-22,Factura X
E-commerce Co,5000.00,Sales Income,Checking,2026-01-22,Depósito venta
\end{verbatim}

\subsection{5. Bank Feeds con Splits (Stripe, PayPal)}

\subsubsection{Escenario real: Depósito de Stripe}

\textbf{Situación:} Stripe deposita \$2,375 que incluye:
\begin{itemize}
    \item Cliente A venda: \$1,500
    \item Cliente B venda: \$900
    \item Stripe fee: -\$25
\end{itemize}

\subsubsection{CSV correctamente formateado}

\begin{verbatim}
bank_feed_date,bank_feed_amount,deposit_id,line_type,
customer_name,amount,account,memo

2026-01-22,2375.00,STRIPE-JAN-001,income,Cliente A,1500.00,
Sales Income,Pago Stripe

2026-01-22,2375.00,STRIPE-JAN-001,income,Cliente B,900.00,
Sales Income,Pago Stripe

2026-01-22,2375.00,STRIPE-JAN-001,fee,,25.00,
Payment Processing Fees,Stripe Fee
\end{verbatim}

\textbf{Validación automática:}
$$1,500 + 900 - 25 = 2,375 \,\checkmark$$

\subsubsection{Cuándo usar bank feeds}

\begin{itemize}
    \item ✅ Pagos en línea (Stripe, PayPal, Mercado Pago)
    \item ✅ Depósitos que agrupan múltiples clientes
    \item ✅ Transacciones con comisiones bancarias
    \item ❌ No para transferencias internas
\end{itemize}

\subsection{6. Reconciliación Bancaria (Mes completo)}

\subsubsection{Flujo paso a paso}

\textbf{Paso 1: Exportar extracto del banco}

En tu banco online:
\begin{itemize}
    \item Seleccionar período (ej: 01-31 enero)
    \item Formato: CSV o Excel
    \item Descargar como: \texttt{banco\_enero.csv}
\end{itemize}

\textbf{Paso 2: Validar formato}

Debe tener columnas:
\begin{verbatim}
date,description,debit,credit,balance,reference
ó
date,description,debit,credit,reference
\end{verbatim}

\textbf{Paso 3: Ejecutar reconciliación}

\begin{verbatim}
Usuario: "Reconcilia enero con banco_enero.csv"

Asistente:
✅ Detectado formato CON balance
✅ Validando matemática:
   Opening: $10,000.00
   Créditos: $8,500.00
   Débitos: $2,300.00
   Ending esperado: $16,200.00
   ✓ Coincide exactamente

✅ Serán creadas transacciones:
   - 5 depósitos (income)
   - 3 bills (gastos)
   - 1 fee bancario

¿Proceder? (sí/no)

Usuario: sí

✅ Reconciliación completada
💾 Reporte: reconciliation_2026-01.csv
\end{verbatim}

\textbf{Paso 4: Validar en QB}

Revisar en QB Online:
\begin{itemize}
    \item Verificar que todas las transacciones aparezcan
    \item Validar que balance de checking coincida
    \item Guardar reporte para auditoría
\end{itemize}

\subsubsection{Formato de ejemplo completo}

\begin{verbatim}
date,description,debit,credit,balance,reference
2026-01-01,Saldo anterior,,,10000.00,OPENING
2026-01-05,Depósito cliente A,,3500.00,13500.00,DEP-001
2026-01-08,Pago servicios,450.00,,13050.00,CHECK-100
2026-01-12,Depósito cliente B,,2200.00,15250.00,DEP-002
2026-01-15,Comisión bancaria,35.00,,15215.00,FEE-01
2026-01-20,Depósito cliente C,,2800.00,18015.00,DEP-003
\end{verbatim}

\subsection{7. Reportes Financieros}

\subsubsection{P\&L (Estado de Resultados)}

\begin{verbatim}
Usuario: "Dame P&L enero-febrero en devengado"

Respuesta: (formato tabla)

ESTADO DE RESULTADOS
Período: 01/01 - 29/02/2026
Método: Devengado

INGRESOS:
  Servicios Profesionales ... $12,500.00
  Ventas de Productos ...... $5,200.00
  ─────────────────────────
  TOTAL INGRESOS ........... $17,700.00

GASTOS:
  Salarios y nómina ....... $6,000.00
  Servicios públicos ...... $500.00
  Alquiler oficina ........ $2,000.00
  Materiales y suministros  $800.00
  ─────────────────────────
  TOTAL GASTOS ............ $9,300.00

GANANCIA NETA ............. $8,400.00

Margen de ganancia: 47.5%
\end{verbatim}

\subsubsection{Balance Sheet (Balance General)}

\begin{verbatim}
Usuario: "Dame balance al 28 de febrero"

BALANCE GENERAL
Fecha: 28/02/2026

ACTIVOS:
  Activos Corrientes:
    Efectivo y bancos ... $18,500.00
    Cuentas por cobrar .. $3,200.00
    Inventario .......... $2,100.00
    Subtotal ............ $23,800.00
  
  Activos Fijos:
    Equipos ............. $5,000.00
    Menos depreciación .. ($500.00)
    Subtotal ............ $4,500.00

  TOTAL ACTIVOS .......... $28,300.00

PASIVOS:
  Pasivos Corrientes:
    Cuentas por pagar ... $2,100.00
    Impuestos por pagar   $1,200.00
    Subtotal ............ $3,300.00

CAPITAL:
  Capital inicial ........ $15,000.00
  Ganancias retenidas .... $10,000.00
  Subtotal .............. $25,000.00

TOTAL PAS + CAP .......... $28,300.00
\end{verbatim}

\subsubsection{Guardar y reutilizar reportes}

\begin{verbatim}
Usuario: "Guarda este P&L como 'Mensual Estándar'"

✅ Guardado con ID: pl_monthly_std_01

Usuario: "Carga 'Mensual Estándar' para marzo"

✅ Mismo reporte, fechas ajustadas (01/03 - 31/03)
\end{verbatim}

\subsection{8. OCR de Facturas PDF}

\subsubsection{Flujo completo: 10 facturas en 2 minutos}

\textbf{Paso 1: Copiar PDFs a Pending bills/}

\begin{verbatim}
Pending bills/
├── factura_amazon_jan.pdf
├── servicio_telefonico.pdf
├── arriendo_inmueble.pdf
├── seguro_empresa.pdf
├── consultoria_auditor.pdf
└── [5 más...]
\end{verbatim}

\textbf{Paso 2: Ejecutar OCR}

\begin{verbatim}
Usuario: "Procesa los bills pendientes"

Asistente:
🔍 Escaneando Pending bills/...
   Encontrados 10 PDFs

🔄 Procesando con OCR...
   📄 factura_amazon_jan.pdf ........... ✓
   📄 servicio_telefonico.pdf .......... ✓
   📄 arriendo_inmueble.pdf ............ ✓
   [+7 más]

✅ Generado: preview_bills_20260122.csv

📊 Resumen OCR:
   - 10 facturas procesadas
   - 9 con confianza >95%
   - 1 con confianza 87% (revisar manualmente)
   - Total: $12,450.80
\end{verbatim}

\textbf{Paso 3: Revisar en Excel}

Abre \texttt{preview\_bills\_20260122.csv} y verifica:

\begin{verbatim}
vendor_name,invoice_number,date,amount,category,description

Amazon,INV-78934,2026-01-18,350.50,Office Supplies,
Papelería variada

Telecom VE,FAC-202601,2026-01-15,145.00,Utilities,
Servicio telefónico enero

Inmobiliario XYZ,RE-2601,2026-01-01,1500.00,Rent,
Arriendo local comercial enero
\end{verbatim}

Corrige si hay errores (vendor, monto, categoría, etc.)

\textbf{Paso 4: Crear bills}

\begin{verbatim}
Usuario: "Aprueba y crea los bills"

Asistente:
✅ Validando CSV...
✅ 10 bills listos

¿Crear en QB? (sí/no)

Usuario: sí

✅ 10 bills creados exitosamente
   - Amazon: $350.50 ✓
   - Telecom VE: $145.00 ✓
   - Inmobiliario: $1,500.00 ✓
   [+7 más]

💾 PDFs movidos a Processed bills/
\end{verbatim}

\subsubsection{Limitaciones y soluciones}

\begin{table}[h]
\centering
\begin{tabular}{|p{3cm}|p{6cm}|}
\hline
\textbf{Problema} & \textbf{Solución} \\
\hline
PDF escaneado & Rescanea con mejor resolución (300+ dpi) \\
\hline
OCR confunde montos & Revisa en Excel, corrige antes de crear \\
\hline
Proveedor no reconocido & Busca manual el vendor ID, úsalo en CSV \\
\hline
Categoría errónea & Cambia en Excel antes de crear bills \\
\hline
PDF corrupto & Cópialo a Processed, digitiza en Excel \\
\hline
\end{tabular}
\caption{OCR troubleshooting}
\end{table}

\newpage

% ============================================
% SECCIÓN 4: FLUJOS DE TRABAJO COMPLETOS
% ============================================
\section{Flujos de Trabajo End-to-End}

\subsection{Flujo 1: Contador autónomo - Mes completo}

\subsubsection{Semana 1-3: Operaciones diarias}

\begin{enumerate}
    \item \textbf{Cada venta:} \texttt{Crea factura a [Cliente] por \$[X]}
    \item \textbf{Cada compra:} \texttt{Crea bill a [Proveedor] por \$[X]}
    \item \textbf{Depósitos diarios:} \texttt{Registra pago de \$[X] de [Cliente]}
    \item \textbf{Fin de semana:} Ejecuta \texttt{¿cuánto he gastado?}
\end{enumerate}

\subsubsection{Semana 4: Cierre de mes}

\begin{enumerate}
    \item Lunes 25: Exportar estado de cuenta del banco
    \item Martes 26: \texttt{Reconcilia el banco con [archivo].csv}
    \item Miércoles 27: Revisar transacciones en QB
    \item Jueves 28: \texttt{Dame P\&L de [mes] en devengado}
    \item Viernes 29: \texttt{Guarda este reporte como 'P\&L [Mes]'}
    \item Sábado 30: \texttt{informe de tokens} (análisis de costos)
\end{enumerate}

\subsubsection{Resultado esperado}

\begin{itemize}
    \item ✅ Contabilidad actualizada (devengado)
    \item ✅ Banco reconciliado
    \item ✅ P\&L generado
    \item ✅ Costos monitoreados
    \item ✅ Tiempo usado: 2-3 horas/mes
\end{itemize}

\subsection{Flujo 2: Firma contable - 20 clientes}

\subsubsection{Sistema de carpetas por cliente}

\begin{verbatim}
clientes/
├── Cliente_A/
│   ├── facturas_enero.csv
│   ├── gastos_enero.csv
│   └── banco_enero.csv
├── Cliente_B/
│   ├── facturas_enero.csv
│   ├── gastos_enero.csv
│   └── banco_enero.csv
└── [18 más...]
\end{verbatim}

\subsubsection{Procesamiento mensual}

\begin{enumerate}
    \item \textbf{Recolección:} Recibir CSVs de todos los clientes
    \item \textbf{Validación:} Revisar que formatos sean correctos
    \item \textbf{Procesamiento:} 
    \begin{verbatim}
    Para cada cliente:
      Procesa facturas_enero.csv
      Procesa gastos_enero.csv
      Reconcilia banco_enero.csv
    \end{verbatim}
    \item \textbf{Reportes:} Generar P\&L para cada cliente
    \item \textbf{Entrega:} Excel con estados financieros
\end{enumerate}

\subsubsection{Tiempo esperado}

\begin{itemize}
    \item Procesamiento: 4-6 horas (vs 20+ horas manual)
    \item Ahorro: 14-16 horas/mes = \$280-320/mes
\end{itemize}

\subsection{Flujo 3: E-commerce - Integración Stripe}

\subsubsection{Configuración mensual}

\begin{enumerate}
    \item Exportar reporte de Stripe: \texttt{CSV con clientes y fees}
    \item Preparar CSV con splits:
    \begin{itemize}
        \item Ventas por cliente (income)
        \item Comisión Stripe (fee)
        \item Validar que sume = depósito bancario
    \end{itemize}
    \item Ejecutar: \texttt{Procesa el bank feed del archivo stripe\_enero.csv}
    \item Resultado: Depósitos creados con cliente y fee correctamente asignados
\end{enumerate}

\subsubsection{Ejemplo: 100 ventas en 1 depósito}

\begin{verbatim}
Stripe deposita: $9,875.50

Desglose (de 100 clientes):
- 100 ventas: $10,000.00
- Stripe fee (1.5%): -$150.00
- Reembolsos: -$50.00
- Ajustes: +$75.50
= Total: $9,875.50 ✓

CSV con 102 líneas (100 income + 1 fee)
Sistema procesa automáticamente
QB recibe 1 depósito con 102 líneas balanceadas
\end{verbatim}

\subsection{Flujo 4: Contabilista - Procesamiento OCR masivo}

\subsubsection{Mensual: 200+ facturas de proveedores}

\begin{enumerate}
    \item \textbf{Recolección:} Stack de 200+ PDFs en Pending bills/
    \item \textbf{OCR en batch:} \texttt{Procesa los bills pendientes}
    \item \textbf{Procesamiento:}
    \begin{itemize}
        \item Sistema escanea todos los PDFs
        \item Genera \texttt{preview\_bills\_[date].csv}
        \item Tiempo: 2-3 minutos
    \end{itemize}
    \item \textbf{Validación:} Excel - revisar errores, corregir
    \item \textbf{Creación:} \texttt{Aprueba y crea los bills}
    \item \textbf{Archivado:} PDFs movidos a Processed bills/
\end{enumerate}

\subsubsection{Ahorro de tiempo}

\begin{itemize}
    \item Digitación manual: 40+ horas/mes
    \item Con OCR + correcciones: 4-5 horas/mes
    \item \textbf{Ahorro: 35 horas/mes = \$700/mes}
\end{itemize}

\newpage

% ============================================
% SECCIÓN 5: FAQ (PREGUNTAS FRECUENTES)
% ============================================
\section{Preguntas Frecuentes (FAQ)}

\subsection{General}

\begin{enumerate}
    \item \textbf{P: ¿El sistema es seguro? ¿Puedo confiar mis datos?}
    
    R: Sí. Los datos se sincronizan directamente con QB Online usando OAuth 2.0. Las credenciales quedan en tu archivo .env local (no se envían a servidores exteriores). El sistema valida transacciones y pide confirmación antes de crear.
    
    \item \textbf{P: ¿Necesito conocimiento técnico para usar TMP AI?}
    
    R: No. El interface es conversacional en español. Solo necesitas saber contabilidad básica (qué es una factura, un bill, depósito). No requiere programación.
    
    \item \textbf{P: ¿Funciona con QB Desktop o solo Online?}
    
    R: Solo QuickBooks Online (QBO). Desktop no tiene API pública. Si usas Desktop, deberías migrar a QBO.
    
    \item \textbf{P: ¿Puedo usar múltiples empresas en QB?}
    
    R: Sí, pero requiere cambiar el REALM\_ID en .env. Para manejar múltiples empresas simultáneamente, necesitarías múltiples instancias con diferentes .env.
    
\end{enumerate}

\subsection{Operaciones}

\begin{enumerate}
    \item \textbf{P: ¿Puedo crear facturas con múltiples líneas diferentes?}
    
    R: El sistema actual soporta líneas simples (1 concepto). Para múltiples líneas diferentes, prepara un CSV especial o créalas manualmente en QB.
    
    \item \textbf{P: ¿Los impuestos se agregan automáticamente?}
    
    R: No. Crea la transacción sin impuesto, luego ajusta manualmente en QB. En futuras versiones (v4.0) habrá soporte automático.
    
    \item \textbf{P: ¿Puedo crear estimaciones (quotes)?}
    
    R: Actualmente no. El sistema crea invoices y bills. Las estimaciones deben crearse en QB.
    
    \item \textbf{P: ¿Soporta multi-moneda?}
    
    R: En v3.0 no. Está planeado para v4.5 (abril 2026). Por ahora, todas las transacciones en moneda de la empresa.
    
\end{enumerate}

\subsection{CSV y Batch Processing}

\begin{enumerate}
    \item \textbf{P: ¿Cuál es el máximo de filas en un CSV?}
    
    R: 500 filas por archivo. Si tienes más, divide en múltiples archivos y procesa uno por uno.
    
    \item \textbf{P: ¿Qué pasa si un cliente no existe en el CSV?}
    
    R: El sistema rechaza esa fila y muestra error. Debes verificar que TODOS los clientes existan en QB antes de procesar.
    
    \item \textbf{P: ¿Puedo revertir un CSV procesado?}
    
    R: No directamente. Debes ir a QB y eliminar transacciones manualmente, o usar QB API para revertir. Siempre haz backup antes de procesar.
    
    \item \textbf{P: ¿El CSV se valida antes de crear?}
    
    R: Sí. El sistema valida:
    \begin{itemize}
        \item Formato de columnas
        \item Tipos de datos (montos, fechas)
        \item Existencia de clientes/proveedores
        \item Disponibilidad de cuentas
    \end{itemize}
    
\end{enumerate}

\subsection{OCR}

\begin{enumerate}
    \item \textbf{P: ¿Qué calidad de PDF necesita el OCR?}
    
    R: Funciona mejor con:
    \begin{itemize}
        \item PDFs nativos (no escaneados)
        \item Texto legible (font > 10pt)
        \item Sin rotaciones o inclinaciones
        \item Resolución mínima 150 dpi
    \end{itemize}
    
    \item \textbf{P: ¿Qué tan preciso es el OCR?}
    
    R: ~95\% en facturas claras. Siempre revisa el CSV generado antes de crear bills. Facturas complejas pueden tener errores en montos o proveedores.
    
    \item \textbf{P: ¿Puedo procesar facturas en otros idiomas?}
    
    R: Sí, pero con menor precisión. El OCR está optimizado para español e inglés.
    
\end{enumerate}

\subsection{Reportes}

\begin{enumerate}
    \item \textbf{P: ¿Cuál es la diferencia entre devengado y caja?}
    
    R: 
    \begin{itemize}
        \item Devengado: Ingresos cuando se factura, gastos cuando se incurren
        \item Caja: Ingresos cuando se recibe dinero, gastos cuando se pagan
    \end{itemize}
    Para contabilidad oficial, usa devengado.
    
    \item \textbf{P: ¿Puedo comparar dos periodos en un reporte?}
    
    R: Actualmente el sistema genera reportes por período. Para comparación, necesitas generar dos reportes y compararlos manualmente en Excel.
    
    \item \textbf{P: ¿Puedo exportar reportes a PDF?}
    
    R: Sí, desde QB Online directamente. El sistema genera tablas en terminal. Cópialo a Word/Excel y exporta a PDF.
    
\end{enumerate}

\subsection{Costos}

\begin{enumerate}
    \item \textbf{P: ¿Cuánto cuesta el sistema?}
    
    R: Depende del consumo de tokens LLM (OpenRouter):
    \begin{itemize}
        \item Operaciones simples: \$0.001-0.01
        \item CSV batch (100 rows): \$0.05-0.10
        \item Reportes: \$0.02-0.05
        \item OCR (5 PDFs): \$0.10-0.20
    \end{itemize}
    Con \$10 de crédito: dura 6.3 meses (optimizado).
    
    \item \textbf{P: ¿Hay costo de QB?}
    
    R: No. El sistema usa QB Online que ya posees (costo separado con Intuit).
    
    \item \textbf{P: ¿Cómo monitoreo mis gastos?}
    
    R: Ejecuta \texttt{¿cuánto he gastado?} para ver sesión actual, o \texttt{informe de tokens} para análisis detallado mensual.
    
\end{enumerate}

\subsection{Troubleshooting General}

\begin{enumerate}
    \item \textbf{P: El programa es muy lento, ¿por qué?}
    
    R: Posibles causas:
    \begin{itemize}
        \item CSV muy grande (>300 rows): divide en archivos más pequeños
        \item Conexión internet lenta: verifica velocidad
        \item QB API sobrecargada: intenta después
        \item LLM (OpenRouter) lento: cambia modelo
    \end{itemize}
    
    \item \textbf{P: ¿Cómo reinicio el programa si se cuelga?}
    
    R: Presiona \texttt{Ctrl+C} para salir, luego \texttt{python3 main.py} para reiniciar.
    
    \item \textbf{P: ¿Hay logs que pueda revisar?}
    
    R: Sí, busca archivo \texttt{tmp\_ai.log} en la carpeta. Contiene detalles de operaciones y errores.
    
\end{enumerate}

\newpage

% ============================================
% SECCIÓN 6: TROUBLESHOOTING AVANZADO
% ============================================
\section{Troubleshooting Avanzado (20+ Errores Comunes)}

\subsection{Errores de autenticación}

\subsubsection{Error 1: "401 Unauthorized - Token expirado"}

\textbf{Síntoma:} \texttt{QB_ACCESS_TOKEN has expired}

\textbf{Causa:} Token de QB vence cada 24 horas.

\textbf{Solución:}
\begin{enumerate}
    \item Sistema intenta refrescar automáticamente usando \texttt{QB_REFRESH_TOKEN}
    \item Si falla, verifica en .env:
    \begin{itemize}
        \item \texttt{QB_REFRESH_TOKEN} es válido
        \item \texttt{QB_CLIENT_ID} y \texttt{QB_CLIENT_SECRET} correctos
        \item No hay caracteres extras (espacios, comillas)
    \end{itemize}
    \item Si sigue fallando, regenera tokens:
    \begin{enumerate}
        \item Ve a developer.intuit.com
        \item Inicia sesión en QB Online
        \item Autoriza la app nuevamente
        \item Copia nuevos tokens a .env
    \end{enumerate}
\end{enumerate}

\subsubsection{Error 2: "Invalid Client ID"}

\textbf{Causa:} Client ID malformado o incorrecto.

\textbf{Solución:}
\begin{enumerate}
    \item Verifica que \texttt{QB_CLIENT_ID} coincida con developer.intuit.com
    \item Comprueba que no tenga espacios al inicio/final
    \item Prueba con una operación simple: \texttt{refrescar chart}
\end{enumerate}

\subsubsection{Error 3: "HTTPS Certificate Error"}

\textbf{Causa:} Problema de SSL/TLS (raro, normalmente en firewalls corporativos).

\textbf{Solución:}
\begin{enumerate}
    \item Intenta desde otra red
    \item Si estás en red corporativa, contacta a IT para permitir openrouter.ai
    \item Como último recurso, desactiva SSL (NO recomendado en producción)
\end{enumerate}

\subsection{Errores de datos}

\subsubsection{Error 4: "Cliente no encontrado"}

\textbf{Síntoma:} \texttt{ACME Corp not found. Did you mean: ACME, ACME International?}

\textbf{Causa:} Nombre exacto no coincide.

\textbf{Solución:}
\begin{enumerate}
    \item Ejecuta: \texttt{Busca cliente ACME}
    \item Elige de la lista sugerida
    \item Si no aparece, cliente no existe en QB. Créalo primero.
    \item Alternativa: usa ID en lugar de nombre: \texttt{Crea factura a ID-123 por \$1000}
\end{enumerate}

\subsubsection{Error 5: "Sumas no cuadran en bank feed"}

\textbf{Síntoma:} \texttt{Deposit DEP-001 = \$100 pero expected \$99.99}

\textbf{Causa:} Error de redondeo o cálculo en CSV.

\textbf{Solución:}
\begin{enumerate}
    \item Abre CSV en Excel
    \item Verifica cada monto con 2 decimales exactos
    \item Usa SUMA() para validar:
    \begin{verbatim}
    =SUM(D2:D100) - E2  (donde E2 es fee)
    Debe = F2 (bank_feed_amount)
    \end{verbatim}
    \item Tolerancia máxima: \$0.01
    \item Si sigue sin cuadrar, hay fila duplicada o valor incorrecto
\end{enumerate}

\subsubsection{Error 6: "CSV formato incorrecto"}

\textbf{Síntoma:} \texttt{Missing required column: 'amount'}

\textbf{Causa:} Encabezados faltantes o mal deletreados.

\textbf{Solución:}
\begin{enumerate}
    \item Ejecuta: \texttt{template csv}
    \item Abre \texttt{depositstemplate.csv} generado
    \item Compara encabezados con tu archivo
    \item Encabezados deben ser exactos (case-sensitive):
    \begin{itemize}
        \item ✅ customer\_name (no Customer Name, no CUSTOMER\_NAME)
        \item ✅ amount (no Amount, Amount., monto)
        \item ✅ from\_account (no de\_account, from account)
    \end{itemize}
\end{enumerate}

\subsubsection{Error 7: "Monto inválido: 'abcd'"}

\textbf{Causa:} Celda con texto en lugar de número.

\textbf{Solución:}
\begin{enumerate}
    \item En Excel, verifica que columnas de montos sean numéricas
    \item Formato correcto: 1500.50 (punto decimal, no coma)
    \item En Latino América: cambiar locales en Excel si usa comas
    \item Busca celdas con texto rojo o símbolos extraños
\end{enumerate}

\subsection{Errores de QB}

\subsubsection{Error 8: "Cuenta no encontrada en QB"}

\textbf{Síntoma:} \texttt{Account 'Checking' not found}

\textbf{Causa:} Nombre de cuenta exacto no coincide con QB.

\textbf{Solución:}
\begin{enumerate}
    \item Ejecuta: \texttt{refrescar chart}
    \item Ejecuta: \texttt{Busca cuenta checking}
    \item Verifica nombre exacto en QB
    \item Usa nombre exacto: "Checking Account" vs "Checking"
\end{enumerate}

\subsubsection{Error 9: "Transacción duplicada"}

\textbf{Síntoma:} Factura para Cliente X existe 2 veces.

\textbf{Causa:} CSV procesado dos veces, o confirmación manual + automática.

\textbf{Solución:}
\begin{enumerate}
    \item Ve a QB Online
    \item Busca transacción duplicada
    \item Elimina la más reciente (o antigua si no la necesitas)
    \item Guarda archivo .csv procesado para no reprocesar
    \item Próxima vez, marca como "procesado" para evitar
\end{enumerate}

\subsubsection{Error 10: "Monto excede límite de QB"}

\textbf{Síntoma:} \texttt{Amount \$1,000,000 exceeds max allowed}

\textbf{Causa:} QB tiene límites por transacción.

\textbf{Solución:}
\begin{enumerate}
    \item Máximo por transacción: \$999,999.99
    \item Divide montos grandes en múltiples transacciones
    \item Ejemplo: \$1,500,000 → 2 depósitos de \$750,000
\end{enumerate}

\subsection{Errores de OCR}

\subsubsection{Error 11: "PDF no reconocido por OCR"}

\textbf{Síntoma:} \texttt{Could not extract text from [factura.pdf]}

\textbf{Causa:} PDF es imagen escaneada de mala calidad.

\textbf{Solución:}
\begin{enumerate}
    \item Rescanea el PDF con mejor resolución (300+ dpi)
    \item Si original es descargado, intenta desde otra fuente
    \item Última opción: digitiza manualmente en Excel
    \item Mueve PDF a Processed bills/ para no reprocesar
\end{enumerate}

\subsubsection{Error 12: "OCR confunde proveedor"}

\textbf{Síntoma:} Detectó "Amazo" en lugar de "Amazon"

\textbf{Causa:} Calidad del PDF, font pequeño o distorsionado.

\textbf{Solución:}
\begin{enumerate}
    \item Revisa el CSV generado: \texttt{preview\_bills\_[date].csv}
    \item Corrige nombre del proveedor en Excel antes de crear
    \item Si proveedor corregido no existe en QB, búscalo y obtén ID
    \item Úsalo en lugar de nombre en CSV
\end{enumerate}

\subsubsection{Error 13: "OCR extrae monto incorrecto"}

\textbf{Síntoma:} Factura por \$1,500 se detectó como \$15.00

\textbf{Causa:} Layout complejo del PDF, múltiples montos en página.

\textbf{Solución:}
\begin{enumerate}
    \item Revisa CSV generado
    \item Corrige monto manualmente en Excel
    \item Verifica que el PDF tenga claramente EL monto a pagar
    \item Para facturas complejas, digitiza manualmente
\end{enumerate}

\subsection{Errores de performance}

\subsubsection{Error 14: "Timeout - operación tardó demasiado"}

\textbf{Síntoma:} \texttt{Operation timed out after 60 seconds}

\textbf{Causa:} CSV muy grande o conexión lenta.

\textbf{Solución:}
\begin{enumerate}
    \item Divide CSV en archivos más pequeños:
    \begin{itemize}
        \item De 500 rows → 2-3 archivos de 150-200 rows
    \end{itemize}
    \item Procesa uno por uno (no todos simultáneamente)
    \item Verifica velocidad de internet: \texttt{speedtest.net}
    \item Si persiste, intenta en horario de menor demanda
\end{enumerate}

\subsubsection{Error 15: "Memoria insuficiente"}

\textbf{Síntoma:} \texttt{MemoryError: Unable to allocate...}

\textbf{Causa:} CSV demasiado grande cargado en memoria.

\textbf{Solución:}
\begin{enumerate}
    \item Máximo recomendado: 500 rows por CSV
    \item Si tienes 1000+, divide en 2+ archivos
    \item Cierra programas innecesarios (navegador, Outlook, etc.)
    \item Reinicia programa entre procesamiento de archivos
\end{enumerate}

\subsection{Errores de conectividad}

\subsubsection{Error 16: "No internet connection"}

\textbf{Síntoma:} \texttt{Connection refused: Cannot reach QB API}

\textbf{Causa:} Sin conexión a internet.

\textbf{Solución:}
\begin{enumerate}
    \item Verifica que WiFi/LAN esté conectado
    \item Ping a google.com desde terminal: \texttt{ping google.com}
    \item Si no responde, conecta a internet
    \item Sistema reintenta automáticamente cada 30 seg
    \item Espera hasta 5 minutos antes de hacer reset
\end{enumerate}

\subsubsection{Error 17: "OpenRouter API unreachable"}

\textbf{Síntoma:} \texttt{Cannot reach openrouter.ai}

\textbf{Causa:} Servicio LLM está caído (raro, 99.9% uptime).

\textbf{Solución:}
\begin{enumerate}
    \item Verifica: \texttt{https://status.openrouter.ai}
    \item Si está caído, espera a que se recupere
    \item Mientras: usa comandos rápidos (\texttt{refrescar chart}, \texttt{template csv})
    \item Estos no requieren LLM
\end{enumerate}

\subsection{Errores de configuración}

\subsubsection{Error 18: ".env file not found"}

\textbf{Causa:} Archivo .env no existe en carpeta principal.

\textbf{Solución:}
\begin{enumerate}
    \item Crea archivo llamado exactamente: \textbf{.env}
    \item Ubícalo en la misma carpeta que main.py
    \item Agrega todas las variables necesarias (ver Sección 2)
    \item Verifica que NO tenga extensión (.env.txt es incorrecto)
\end{enumerate}

\subsubsection{Error 19: "QB_REALM_ID invalid format"}

\textbf{Causa:} Realm ID contiene caracteres inválidos.

\textbf{Solución:}
\begin{enumerate}
    \item Realm ID debe ser números solamente: 1234567890
    \item Obtén en developer.intuit.com → Apps → Settings
    \item Verifica que no tenga guiones, letras u otros caracteres
\end{enumerate}

\subsubsection{Error 20: "OPENROUTER_API_KEY invalid"}

\textbf{Causa:} API key expirada o incorrecta.

\textbf{Solución:}
\begin{enumerate}
    \item Ve a openrouter.ai
    \item Inicia sesión
    \item Regenera API key en Settings
    \item Copia el valor completo: \texttt{sk-or-v1-xxxxxxxx}
    \item Reemplaza en .env
\end{enumerate}

\newpage

% ============================================
% SECCIÓN 7: MÉTRICAS Y MONITOREO
% ============================================
\section{Métricas y Monitoreo}

\subsection{Comando: ¿cuánto he gastado?}

\subsubsection{Salida típica}

\begin{verbatim}
═══════════════════════════════════════════════
    RESUMEN DE SESIÓN - Enero 22, 2026
═══════════════════════════════════════════════

Duración: 1 hora 25 minutos
Inicio: 10:30 AM | Fin: 11:55 AM

TOKENS USADOS:
  Input tokens: 45,230
  Output tokens: 12,100
  Total: 57,330 tokens

COSTOS:
  Costo estimado: $0.23
  Presupuesto restante: $9.77

OPERACIONES REALIZADAS:
  ✓ Búsquedas: 5
  ✓ Facturas creadas: 3
  ✓ Bills creados: 2
  ✓ Depósitos: 1
  ✓ Reportes: 1
  ────────────────
  Total operaciones: 12

EFICIENCIA:
  Tokens por operación: 4,778
  Costo por operación: $0.019
  Operaciones por hora: 8.4

═══════════════════════════════════════════════
\end{verbatim}

\subsection{Comando: informe de tokens}

Genera archivo \textbf{tokenusagereport.xlsx} con:

\subsubsection{Hoja 1: Resumen Mensual}

\begin{table}[h]
\centering
\begin{tabular}{|l|r|r|r|}
\hline
\textbf{Mes} & \textbf{Sesiones} & \textbf{Tokens} & \textbf{Costo USD} \\
\hline
Enero 2026 & 12 & 687,960 & \$2.75 \\
Diciembre 2025 & 8 & 450,120 & \$1.80 \\
Noviembre 2025 & 15 & 892,340 & \$3.57 \\
\hline
\textbf{TOTAL} & \textbf{35} & \textbf{2,030,420} & \textbf{\$8.12} \\
\hline
\end{tabular}
\caption{Uso de tokens por mes}
\end{table}

\subsubsection{Hoja 2: Análisis por Tipo de Operación}

\begin{table}[h]
\centering
\begin{tabular}{|l|r|r|r|}
\hline
\textbf{Operación} & \textbf{Frecuencia} & \textbf{Tokens Totales} & \textbf{\% del Total} \\
\hline
Búsquedas simples & 45 & 90,000 & 4.4\% \\
Crear facturas & 28 & 240,000 & 11.8\% \\
Crear bills & 22 & 190,000 & 9.4\% \\
CSV batch (depósitos) & 8 & 560,000 & 27.6\% \\
OCR facturas & 12 & 680,000 & 33.5\% \\
Reportes & 15 & 270,000 & 13.3\% \\
\hline
\end{tabular}
\caption{Tokens por tipo de operación}
\end{table}

\subsubsection{Insights y recomendaciones}

\begin{itemize}
    \item ✅ OCR es la operación más cara (33.5\%) pero más productiva
    \item ✅ CSV batch es 40\% más eficiente que operaciones individuales
    \item ⚠️ Reportes pueden optimizarse (guardar y reutilizar reduce 30\%)
    \item 💡 Usar comandos rápidos ahorra 5-10\% de costo/sesión
\end{itemize}

\subsection{Proyecciones de presupuesto}

\textbf{Crédito: \$10}

\subsubsection{Escenario 1: Sin optimizaciones}

\begin{itemize}
    \item Operaciones ineficientes
    \item OCR por PDF individual
    \item Reportes sin reutilización
    \item \textbf{Duración: 2.7 meses}
\end{itemize}

\subsubsection{Escenario 2: Con optimizaciones (recomendado)}

\begin{itemize}
    \item Batch processing
    \item Reportes guardados y reutilizados
    \item Comandos rápidos para operaciones simples
    \item OCR masivo (5+ PDFs juntos)
    \item \textbf{Duración: 6.3 meses} → \$1.59/mes
\end{itemize}

\subsubsection{Escenario 3: Alta productividad}

\begin{itemize}
    \item 100+ transacciones batch/semana
    \item Reportes reutilizados constantemente
    \item Máximo uso de comandos rápidos
    \item \textbf{Duración: 8+ meses} → <\$1.25/mes
\end{itemize}

\newpage

% ============================================
% SECCIÓN 8: INTEGRACIONES Y EXTENSIONES
% ============================================
\section{Integraciones con Herramientas Externas}

\subsection{Integración: Stripe}

\subsubsection{Paso 1: Exportar reporte de Stripe}

\begin{enumerate}
    \item Stripe Dashboard → Balances → Transactions
    \item Seleccionar período (ej: 01-31 enero)
    \item Download CSV
    \item Archivo contiene: date, amount, customers, fees
\end{enumerate}

\subsubsection{Paso 2: Transformar a formato compatible}

El CSV de Stripe NO es directamente compatible. Necesita transformación:

\begin{verbatim}
Stripe CSV original:
date,customer_name,amount,stripe_fee,net

Transformar a:
bank_feed_date,bank_feed_amount,deposit_id,line_type,
customer_name,amount,account,memo
\end{verbatim}

\subsubsection{Paso 3: Procesar con TMP AI}

\begin{verbatim}
Usuario: "Procesa el bank feed stripe_enero.csv"

Sistema:
✅ Valida que sumas cuadren
✅ Agrupa por deposit_id
✅ Crea un único depósito por día con splits
\end{verbatim}

\subsection{Integración: PayPal}

Similar a Stripe:

\begin{enumerate}
    \item PayPal → Activity → Download CSV
    \item Transformar formato (date, customer, amount, fee)
    \item Procesar con TMP AI banco feed
\end{enumerate}

\subsection{Integración: Google Sheets}

\subsubsection{Exportar datos de QB a Sheets}

\begin{enumerate}
    \item Generar P\&L o reporte en TMP AI
    \item Copiar tabla de terminal
    \item Pegar en Google Sheets
    \item Formatear y compartir con stakeholders
\end{enumerate}

\subsubsection{Usar Sheets como entrada}

\begin{enumerate}
    \item Llenar datos en Google Sheets (compartido con equipo)
    \item Descargar como CSV: \texttt{File → Download → CSV}
    \item Procesar con TMP AI
\end{enumerate}

\subsection{Integración: Excel (análisis posterior)}

\begin{enumerate}
    \item Generar reporte con TMP AI
    \item Copiar salida a Excel
    \item Crear tablas dinámicas (pivot tables)
    \item Agregar gráficos para presentaciones
    \item Compartir con clientes/stakeholders
\end{enumerate}

\newpage

% ============================================
% SECCIÓN 9: ROADMAP Y FUTURO
% ============================================
\section{Roadmap de Desarrollo}

\subsection{v3.0 (Actual - Enero 2026)}

✅ 31 funciones core
✅ OCR de facturas
✅ Bank feed processing
✅ Optimización de tokens 57\%
✅ Reconciliación bancaria

\subsection{v3.5 (Febrero 2026)}

\begin{itemize}
    \item ✓ Soporte para impuestos automáticos
    \item ✓ Mejoras en fuzzy matching
    \item ✓ Multi-idioma (inglés, portugués)
    \item ✓ Auditoria (log de quién hizo qué)
\end{itemize}

\subsection{v4.0 (Marzo 2026)}

\begin{itemize}
    \item ✓ Integración nativa Stripe API
    \item ✓ Integración nativa PayPal API
    \item ✓ Depósitos automáticos diarios
    \item ✓ Alertas de anomalías
\end{itemize}

\subsection{v4.5 (Abril 2026)}

\begin{itemize}
    \item ✓ Multi-moneda automático
    \item ✓ Soporte de descuentos complejos
    \item ✓ Estimaciones (quotes)
    \item ✓ Reportes comparativos
\end{itemize}

\subsection{v5.0 (Junio 2026)}

\begin{itemize}
    \item ✓ Asientos de diario automáticos
    \item ✓ Depreciación de activos
    \item ✓ Integración con Zapier
    \item ✓ Dashboard web interactivo
    \item ✓ API pública para terceros
\end{itemize}

\newpage

% ============================================
% SECCIÓN 10: SOPORTE Y CONTACTO
% ============================================
\section{Soporte y Contacto}

\subsection{Documentación}

\begin{enumerate}
    \item \textbf{CONTEXT.md} - Documentación técnica completa
    \item \textbf{main.py} - Código fuente comentado (95K+ caracteres)
    \item \textbf{Este manual} - Guía usuario final
    \item \textbf{FAQ} - Preguntas frecuentes (50+ tópicos)
    \item \textbf{Logs} - Ver \texttt{tmp\_ai.log} para debug
\end{enumerate}

\subsection{Reportar errores}

\textbf{Información a incluir:}

\begin{enumerate}
    \item Descripción exacta de qué hiciste
    \item Mensaje de error completo (copiar-pegar)
    \item Archivo que causó error (sin datos sensibles)
    \item Versión Python: \texttt{python3 --version}
    \item Sistema operativo: Windows / Mac / Linux
    \item ¿Estás en sandbox o producción?
\end{enumerate}

\subsection{Sugerencias de mejora}

Todas las sugerencias son bienvenidas:

\begin{itemize}
    \item Nuevas operaciones que necesites
    \item Mejoras de interfaz o claridad
    \item Optimizaciones de performance
    \item Integraciones adicionales
\end{itemize}

\subsection{Información de versión}

\begin{tcolorbox}[colback=lightblue, colframe=darkblue, boxrule=2pt]
    \begin{itemize}
        \item \textbf{Sistema:} TMP AI Assistant v3.0
        \item \textbf{Fecha:} Enero 22, 2026
        \item \textbf{Desarrollador:} Alfredo
        \item \textbf{LLM Backend:} DeepSeek V3 vía OpenRouter
        \item \textbf{QB API:} v4.0
        \item \textbf{Python:} 3.9+
        \item \textbf{Estado:} Producción
        \item \textbf{Uptime SLA:} 99.5\% (sin garantía)
    \end{itemize}
\end{tcolorbox}

\newpage

% ============================================
% APÉNDICES
% ============================================
\appendix

\section{Checklist: Primeros 7 días}

\begin{itemize}
    \item [ ] \textbf{Día 1:} Instalar Python, crear .env, ejecutar main.py
    \item [ ] \textbf{Día 2:} Ejecutar búsqueda de cliente test, refrescar chart
    \item [ ] \textbf{Día 3:} Crear factura de prueba (\$100)
    \item [ ] \textbf{Día 4:} Crear bill de prueba (\$50)
    \item [ ] \textbf{Día 5:} Procesar CSV pequeño (2-3 depósitos)
    \item [ ] \textbf{Día 6:} Generar P\&L
    \item [ ] \textbf{Día 7:} Ejecutar \texttt{informe de tokens}
\end{itemize}

\section{Plantillas de Comandos (Copy-Paste)}

\subsection{Búsquedas}
\begin{verbatim}
Busca cliente NOMBRE_PARCIAL
Busca proveedor NOMBRE
Busca cuenta ACCOUNT_NAME
\end{verbatim}

\subsection{Transacciones}
\begin{verbatim}
Crea factura a CLIENTE por $MONTO por CONCEPTO
Crea bill a PROVEEDOR por $MONTO por CONCEPTO
Registra pago de $MONTO de CLIENTE a CUENTA
\end{verbatim}

\subsection{Batch}
\begin{verbatim}
Procesa archivo_deposits.csv
Procesa el bank feed archivo_stripe.csv
Reconcilia el banco con estado_cuenta.csv
\end{verbatim}

\subsection{Reportes}
\begin{verbatim}
Dame P&L de ENERO en devengado
Genera balance al 31 de enero
Guarda este reporte como "NOMBRE"
Carga el reporte "NOMBRE"
\end{verbatim}

\subsection{OCR}
\begin{verbatim}
Procesa los bills pendientes
Aprueba y crea los bills del CSV
\end{verbatim}

\subsection{Mantenimiento}
\begin{verbatim}
¿cuánto he gastado?
informe de tokens
refrescar chart
template csv
listar reportes
\end{verbatim}

\section{Glosario de Términos}

\begin{itemize}
    \item \textbf{QB / QBO:} QuickBooks Online
    \item \textbf{Invoice:} Factura (documento de venta)
    \item \textbf{Bill:} Cuenta por pagar (gasto)
    \item \textbf{Deposit:} Depósito (movimiento de dinero)
    \item \textbf{Chart of Accounts:} Catálogo de cuentas
    \item \textbf{P\&L:} Estado de resultados (Profit \& Loss)
    \item \textbf{Bank Feed:} Depósito con múltiples líneas
    \item \textbf{CSV:} Archivo con datos tabulados
    \item \textbf{OCR:} Escaneo inteligente de PDFs
    \item \textbf{Token:} Unidad de costo de API LLM
    \item \textbf{Devengado:} Ingresos cuando se facturan, gastos cuando se incurren
    \item \textbf{Caja:} Ingresos cuando se reciben, gastos cuando se pagan
    \item \textbf{Reconciliación:} Validación de que QB coincida con banco
    \item \textbf{Fuzzy Matching:} Búsqueda aproximada (no requiere exactitud)
    \item \textbf{API:} Interfaz para conectar aplicaciones
    \item \textbf{OAuth 2.0:} Protocolo seguro de autenticación
\end{itemize}

\section{Índice de Errores Comunes}

\begin{enumerate}
    \item 401 Unauthorized → Token expirado
    \item Invalid Client ID → Cliente mal configurado
    \item Cliente no encontrado → Nombre inexacto
    \item Sumas no cuadran → Error en CSV
    \item CSV formato incorrecto → Encabezados faltantes
    \item Monto inválido → Celda contiene texto
    \item Cuenta no encontrada → Nombre no existe
    \item Transacción duplicada → CSV procesado dos veces
    \item Monto excede límite → Máximo \$999,999.99
    \item PDF no reconocido → OCR calidad baja
    \item OCR confunde datos → Revisar y corregir CSV
    \item Timeout → CSV demasiado grande
    \item No internet → Verificar conexión
    \item .env not found → Crear archivo en carpeta principal
    \item API key inválida → Regenerar en OpenRouter
\end{enumerate}

\newpage

\begin{center}
    \vspace{3cm}
    \Large\textbf{\textcolor{darkblue}{¿LISTO PARA AUTOMATIZAR?}}
    
    \vspace{1cm}
    \texttt{python3 main.py}
    
    \vspace{2cm}
    \Large\textcolor{accent}{TMP AI Assistant v3.0}
    
    \vspace{0.5cm}
    \large Automatización Inteligente de Contabilidad en QuickBooks
    
    \vspace{3cm}
    \small\textcolor{darkblue}{Hecho con \& por Alfredo | Enero 2026}
\end{center}

\end{document}